% Compile from report/ after running `python run.py report` at the repository root.
% The whole report, including outside sources, is kept within four pages.
\documentclass[11pt,a4paper]{article}
\usepackage[margin=2.2cm]{geometry}
\usepackage{graphicx,booktabs,hyperref}
\usepackage{lmodern}
\usepackage{float}
\hypersetup{hidelinks,
  pdftitle={SSU Data Dash 2026: Analysts},
  pdfauthor={Farhan Zahid, Haojie Wang, Jiaming Han, Luke Zhenyang He}}
\graphicspath{{../figures/}}
\setlength{\parskip}{2pt}
\setlength{\parindent}{0pt}
\setlength{\emergencystretch}{2em}
\setlength{\textfloatsep}{8pt}
\setlength{\intextsep}{6pt}
\setlength{\abovecaptionskip}{4pt}
% Keep the template's 11pt text and section style, with less heading whitespace.
\makeatletter
\renewcommand\section{\@startsection{section}{1}{\z@}
  {-2.4ex \@plus -0.8ex \@minus -0.2ex}
  {1.2ex \@plus 0.2ex}{\normalfont\Large\bfseries}}
\makeatother

\title{SSU Data Dash 2026: Analysts}
% Alphabetical by full English name.
\author{Farhan Zahid \and Haojie Wang \and Jiaming Han \and Luke Zhenyang He}
\date{}

\begin{document}
\maketitle
\vspace{-1.7em}

\noindent\begin{tabular}{@{}lp{0.64\linewidth}@{}}
\textbf{Video (YouTube):} & \textit{[https://youtu.be/BEEqbytFvJA]} \\
\textbf{Code and submission files:} & \url{https://github.com/LucasWHJ/SSU-Data-Dash} \\
\end{tabular}

\section*{Summary}
We built separate models to value homes with and without a list price, and predicted their sale prices.
Histogram-based gradient-boosted trees achieved validation MdAPE of 7.76\% off-market and 4.10\% list-aware on
10,069 homes per track. We found that list price provided a useful anchor, while off-market MdAPE varied from
6.5\% in the lowest predicted-price decile to 12.2\% in the highest.
Validation scores may still flatter test performance, as selection was repeated on the same six months,
and freezing training data at August 2025 latest raised MdAPE to 8.39\% off-market and 4.31\% list-aware.

\section{Data and features}
We used the organisers' 150,000 sales synthetic dataset, covering homes with sold dates from January 2021 to May 2026
(\href{https://huggingface.co/datasets/Baliyang/SSU_data_dash}{competition dataset}).
The hidden test files contained 3,974 off-market and 4,064 list-aware homes. Sold price supplied the
training target, and sold date defined historical splits, neither entering the prediction features in our final model.
We also excluded days on market and the external automated valuation model (AVM) estimate, as both were
absent from the test files.
For final predictions, both models were refitted on all 150,000 known sales.

Both tracks used features such house-area midpoints (in square feet), num bedrooms, num bathrooms, parking capacity, garage capacity, lot dimensions
and lot area as numerical features, and house type, style, municipality and community as unchanged categorical features. Trees handled missing numerical values,
while categories were learned from each training fold and unseen levels were mapped to missing.
All time features came from the listing date (\verb|date_listed|). These included months since
January 2021 as a linear progression, and either calendar month/quarter, or sin/cos month harmonics as wrap-around features.

We used descriptions to capture details missing from the structured columns. We first used SQL
\texttt{ILIKE} queries on the \texttt{remarks} column across \verb|listings|,
\verb|test_off_market| and \verb|test_list_aware| to count the number of listings matching each candidate pattern.
This frequency screening identified eleven candidate feature groups for validation: renovation,
repair needs, move-in readiness, transit, ravine, views, basement, balcony/terrace, locker/storage
unit, original finishes, and full renovation. We then encoded each group as a 0/1 indicator using
case-insensitive regex matching in Python, and also included description word count as a
numerical feature.

The original-finish and full-renovation indicators matched the word \textit{original} and phrases
such as \textit{fully redone} and \textit{top to bottom}, respectively. They appeared in 24.85\% and 13.61\%
of the 139,931 sales before validation began. We kept \textit{original} separate from repair needs
because original finishes could still be well maintained.

A tree splits one coordinate at a time, so uncovering diagonal geographic patterns can require many inefficient splits.
Off-market used ten fixed rotations of latitude and scaled longitude at 15--165 degrees,
excluding 90 degrees because that projection would reproduce scaled longitude.
It also used area and room ratios alongside actual area size and room counts. List-aware kept raw
coordinates and sin/cos month encoding, plus asking price and asking price per square foot.
Since we selected each track's full configuration on validation, these results do not isolate
the list-aware effects of rotation or seasonal encoding.

\newpage
\section{Models}
\textbf{Off-market.} A \texttt{HistGradientBoostingRegressor} learned log sold price, and we
exponentiated its predictions (\verb|np.exp()|) to obtain dollar estimates. Without a list price anchor, the model had to learn the price level.
Squared error on log prices gave more weight to large proportional misses while reducing the
dominance of expensive homes' dollar errors. Settings were learning rate 0.18, 400 iterations,
60 maximum leaves, 30 minimum samples per leaf and 70\% feature sampling at each split.

\textbf{List-aware.} A separate \texttt{HistGradientBoostingRegressor} learned $\log(\mathrm{sold}/\mathrm{list})$, then
exponentiated its output and re-multiplied by list price. Absolute error limited the influence of
unusually large departures from this anchor and targeted the conditional median log ratio.
Exponentiation preserves that median. Settings were learning rate 0.06,
260 iterations, 31 leaves, 40 samples per leaf and all features at each split. Both tracks used L2
regularisation 0.1, seed 0, and no early stopping. Since MdAPE summarises percentage errors across
homes, we selected whole configurations through validation rather than assuming that either training
loss optimised this metric.

\section{Validation}
For each month from December 2025 through May 2026, we fitted a fresh model using only sales
completed before that month, then predicted every home sold during that month. Training grew from
139,931 to 147,709 sales, while the six months contributed 10,069 distinct validation homes.
Earlier validation months entered subsequent training sets, matching the starter's expanding-window
procedure, and metrics pooled all six months.

Prediction inputs excluded post-sale fields and off-market asking price. Features used each home's
available attributes, with training-fold categorical schemas. No hidden-test prices or future price
averages were used. Splits followed sale months, not each home's original listing date.

For sold price $y$ and prediction $\hat y$, we computed signed error as $100(\hat y-y)/y$,
from which we obtained MdAPE and MAPE by taking the median and mean absolute values.
Positive median signed error indicates upward bias in house valuation.
Bootstrap 95\% intervals used 1,000 home resamples (seed 0), excluding model-selection uncertainty
and dependence from shared monthly shocks.

\section{Results}
\begin{table}[ht]
\centering
\resizebox{\linewidth}{!}{\input{../reports/results_table.tex}}
\caption{Walk-forward validation on the listings. MdAPE: median absolute percentage error.
MAPE: mean absolute percentage error. Brackets: bootstrap 95\% interval.}
\end{table}

Our MdAPE was lower by 13.04\% points off-market and 3.20\% list-aware than the baseline
reported starter scores of 20.8\% (nearest sales) and 7.3\% (asking price). Small overall median
signed errors, +0.35\% and +0.31\%, motivated a closer calibration check.

After frequency screening, we compared the first nine groups in Section 1 against all eleven,
keeping other features, model parameters and validation months fixed. Adding \textit{original} and
full-renovation indicators lowered off-market MdAPE by 0.11\% points and list-aware by
only 0.01\%, but paired bootstrap intervals for both changes included zero. We retained all eleven
for their lower observed errors, while treating the gain as uncertain.

\begin{table}[H]
\centering
\small
\begin{tabular}{p{8.0cm}rr}
\toprule
Description features & Off-market MdAPE & List-aware MdAPE \\
\midrule
Nine keyword groups and word count & 7.87\% & 4.11\% \\
Eleven keyword groups and word count (selected) & 7.76\% & 4.10\% \\
\bottomrule
\end{tabular}
\caption{Keyword comparison with other features and model settings held fixed.}
\end{table}

\begin{samepage}
We expected sentence embeddings and TF-IDF weights to capture more than keywords did, but neither
improved earlier-model MdAPE. Qwen condition extraction was evaluated only on a small sample.
\par
\end{samepage}

\begin{samepage}
An off-market notebook trial also added 1/30/90-day rolling means for log sale prices by house type,
municipality and house style, using sales before listing and no later than the training cutoff.
MdAPE/MAPE results were 8.03\%/10.67\% compared to 7.86\%/10.49\% for our earlier configuration, so these
features were not included. A sale-date-based variant (using sales before a home's sold date instead of list date) was
compared as well, with no improvements made to MdAPE or MAPE. This was done as a diagnostic only,
because sale dates are meant to be unknown at prediction time.
\par
\end{samepage}

\section{Where the model fails}
\begin{figure}[H]
\centering
\includegraphics[width=0.72\linewidth]{error_by_price_range.png}
\caption{Absolute percentage error grouped by actual sale price. Lines show medians and bands show
the 25th--75th percentiles. The predicted-price check below also found larger high-end errors.}
\end{figure}

\textbf{Price and calibration.} Grouping by actual sale price initially suggested that we were
overvaluing cheaper homes and undervaluing expensive ones. We first interpreted this as shrinkage toward
typical prices, although grouping on the realised outcome can itself create that pattern.
Regrouping by predicted price (Table 3) left median overestimation of 3.01\% at the low end and
0.30\% at the high end. The high end negative bias largely disappeared, while low end overestimation
remained. The list-aware track's corresponding errors were closer to zero.

Typical errors were still larger for expensive homes. By predicted price, off-market MdAPE was 6.46\% in the lowest
decile and 12.23\% in the highest, whereas list-aware values were 3.39\% and 4.76\%.
The overall 90th-percentile absolute error was 21.68\% off-market and 11.05\% list-aware.

\begin{table}[ht]
\centering
\small
\begin{tabular}{llrr}
\toprule
Track & Deciles defined by & Lowest: signed error & Highest: signed error \\
\midrule
Off-market & Actual sale price & +7.24\% & $-$7.31\% \\
Off-market & Predicted price & +3.01\% & +0.30\% \\
List-aware & Actual sale price & +2.05\% & $-$1.47\% \\
List-aware & Predicted price & +0.90\% & +0.22\% \\
\bottomrule
\end{tabular}
\caption{Median signed errors. Membership changes when grouping by predicted rather than actual price.}
\end{table}

\textbf{Property type.} Off-market MdAPE was 6.65\% for condo apartments and 9.53\% for detached
homes. Co-ops showed 12.19\% off-market versus 4.93\% list-aware, but only 48 were validated,
so this contrast needs caution.

\begin{samepage}
\textbf{Time and neighbourhood.} Monthly MdAPE ranged from 7.46\% to 8.18\% off-market and
3.81\% to 4.35\% list-aware, which peaked in May. Elmview had the highest off-market MdAPE among
communities with at least 100 validation homes (10.87\%), although its 184 homes represented only
1.8\% of the validation sample.
\par
\end{samepage}

\section{Limitations}
\textbf{Selection uncertainty.} Repeated selection can overfit these six months, and the 0.11- and
0.01-point keyword gains are within bootstrap uncertainty. With off-market feature sampling
(\texttt{max\_features = 0.7}), reordering columns can change the randomly selected candidate
features even with a fixed random seed. Our bootstrap intervals do not capture variation from column
order or random seed, being based off our predictions only. Hidden-test accuracy is also unknown, and we would
confirm choices on reserved future months.

\textbf{Test-description screening.} Keyword selection used unlabelled descriptions from both
test files, although their sale prices remained hidden.

\textbf{Prediction timing.} Of 10,069 validation homes, 6,465 (64.2\%) were listed before their
validation month's first day, which was the training cutoff. Training could therefore use sales
unavailable at listing, making these backtests potentially optimistic for listing-time valuation.
Valuing at listing time would need a separate backtest with a cutoff for each listing.

\textbf{Coverage and horizon.} The test period extends further past training than validation did.
In validation, 60.7\% of listing dates fell inside each fold's observed date range. This is below
64.2\% because the latest training listing preceded the month-start cutoff. The furthest date was
32 days beyond the range. In the test files, 71.4\% of off-market and 72.4\% of list-aware dates
exceeded the final training maximum, by up to 122 days. For the linear trend input, later values
enter the same extreme bin, so this feature cannot continue a rising or falling market.

\textbf{Horizon stress test.} We froze training at August 31, 2025 (133,779 sales) and evaluated the
same December--May homes. MdAPE rose from 7.76\% to 8.39\% off-market and from 4.10\% to 4.31\%
list-aware. This combines a longer horizon with less recent training data, so it is a stress test
rather than a forecast of hidden-test error. Removing the linear trend in the same frozen setup
raised errors to 14.40\% and 7.10\%. We have not tested an extrapolating time-based model.
The smaller list-aware increase is consistent with some benefit from the asking-price
anchor, although both tracks worsened substantially when the trend was removed.

\textbf{Data and drift.} As the data is synthetic, area midpoints lose information, and descriptions
may omit or overstate details, so real-world use would need separate validation. List-aware's
highest monthly error came in May, but a single month cannot establish drift.

\textbf{Decision risk.} Off-market misses sometimes exceeded 20\%. We would check recent
comparables and estimate prediction intervals for individual homes.

\section*{Outside data}
Final predictions used only competition data. The exploratory models below used pre-June 2026
revisions, although their pretraining periods were unverified.
Their outputs did not enter final features or predictions, and no real supplementary listing or sales records were used.

\begingroup
\centering
\small
\begin{tabular}{p{3.0cm}p{7.3cm}p{3.3cm}}
\toprule
Source & URL & Periods covered \\
\midrule
MiniLM &
\href{https://huggingface.co/sentence-transformers/all-MiniLM-L6-v2/tree/c9745ed1d9f207416be6d2e6f8de32d1f16199bf}{\nolinkurl{https://huggingface.co/sentence-transformers/all-MiniLM-L6-v2}} &
Revision before June 2026 \\
Qwen2.5 1.5B &
\href{https://huggingface.co/mlx-community/Qwen2.5-1.5B-Instruct-4bit/tree/8b403126fc14f14cfc99bb4cfa72ecbc129ea677}{\nolinkurl{https://huggingface.co/mlx-community/Qwen2.5-1.5B-Instruct-4bit}} &
Revision before June 2026 \\
\bottomrule
\end{tabular}
\par
\endgroup

\section*{Reproducing our results}
With Python 3.11, run from the repository root (set \verb|DATA_DIR| only for local data):
\begin{verbatim}
python -m pip install -r requirements.txt && python run.py validate
python run.py report && python run.py predict && python run.py check
\end{verbatim}
\end{document}
